\chapter{Chiral Hochschild calculus on $\mathbf H_{\Delta_5}$ at
$c=-214$}
\label{chap:hochschild-calculus}

The obstruction classes $[e_k]\in\ChirHoch^3(\mathbf H_{\Delta_5})$
are represented at chain level by Virasoro quasi-primary composites
under the $\CohoZ_3$-projection. The K$3\times E$ CoHA pairing
$\langle[\chi_k],[e_k]\rangle_{\Phi_k}$ reads off specific rationals
at the universal central charge $c=-214$ attached to the
$\Delta_5$-chamber of the $\mathrm K3$-Yangian moduli.

\section{Weight-$4$ cocycle}
\label{sec:e4-calculus}

\begin{theorem}[$e_4$ chain-level form]
\label{thm:e4-vol3}
\ClaimStatusProvedHere
On the Virasoro subalgebra $\Vir_{c=-214}\subset\mathbf H_{\Delta_5}$,
\begin{equation}
\label{eq:e4-vol3}
e_4(z)={:}T\partial^2T{:}-\tfrac{3}{2}{:}(\partial T)^2{:}
     +\tfrac{321}{10}\partial^4T+\hbar\,\mathrm{qt}(J^{(4)}),
\end{equation}
with universal-$c$ coefficients
$\alpha_4^{(1)}=1,\alpha_4^{(2)}=-3/2,\alpha_4^{(3)}=-3c/20$.
\end{theorem}

\begin{proof}
Using $L_1T=0$ and $L_1\partial^nT=n(n+1)\partial^{n-1}T$,
\begin{align*}
L_1{:}T\partial^2T{:}&=6{:}T\partial T{:}+3c\,\partial^3T,\\
L_1{:}(\partial T)^2{:}&=4{:}T\partial T{:},\\
L_1\partial^4T&=20\partial^3T.
\end{align*}
Annihilating the ${:}T\partial T{:}$-coefficient gives
$6\alpha_4^{(1)}+4\alpha_4^{(2)}=0$, hence
$\alpha_4^{(2)}=-3/2\cdot\alpha_4^{(1)}$; annihilating
$\partial^3T$ gives
$3c+20\alpha_4^{(3)}=0$, hence
$\alpha_4^{(3)}=-3c/20$. At $c=-214$ this is $321/10$.
Primary: Zamolodchikov~$1985$~TMF~$65$;
Bais--Bouwknegt--Schoutens--Surridge~$1988$~NPB~$304$.
\end{proof}

\begin{theorem}[CoHA-Casimir pairing at weight $4$]
\label{thm:chi4-e4-vol3}
\ClaimStatusProvedHere
At $c=-214$, on K$3\times E$,
\begin{equation}
\label{eq:chi4-e4-pairing-vol3}
\langle[\chi_4],[e_4]\rangle_{\Phi_4}
=\tfrac{65193}{10}\,\mathrm{Vol}(E)\,(2\pi\mathrm i)^4.
\end{equation}
\end{theorem}

\begin{proof}
Schiffmann--Vasserot~$2013$~\S$6$.$3$ CoHA-Casimir
$\mathrm{Cas}_4(\alpha)=\chi(\mathcal O_{\mathrm K3})
\bigl(1+c(c+11)/20\bigr)\,\mathrm{Vol}(E)(2\pi\mathrm i)^4$
with $\chi(\mathcal O_{\mathrm K3})=2$ and
$1+(-214)(-203)/20=21731/10$; the pairing
$\langle[\chi_4],[e_4]\rangle_{\Phi_4}=\tfrac{3}{2}\mathrm{Cas}_4
=\tfrac{3}{2}\cdot 2\cdot 21731/10\cdot\mathrm{Vol}(E)(2\pi\mathrm i)^4
=65193/10\cdot\mathrm{Vol}(E)(2\pi\mathrm i)^4$.
Cross-verification: Oberdieck~$2018$~Geom.\ \& Topol.~$22$ Thm.~$2$
gives $[p^{-1}q^{-1}y^4](-\Phi_{10}^{-1})=21731/10$ in the
Gritsenko--Nikulin positive Weyl chamber, matching path (A).
CY-$2$ formality (Kontsevich~$2003$) forces linear-in-$\mathrm{Vol}(E)$
scaling, matching the leading behaviour.
\end{proof}

\section{Weight-$5$ cocycle}
\label{sec:e5-calculus}

\begin{theorem}[$e_5$ chain-level form]
\label{thm:e5-vol3}
\ClaimStatusProvedHere
The weight-$5$ quasi-primary obstruction cocycle is
\begin{equation}
\label{eq:e5-vol3}
e_5(z)={:}TT\partial T{:}-\tfrac{3}{10}{:}T\partial^3T{:}
     +\tfrac{107}{28}{:}(\partial T)(\partial^2T){:}
     -\tfrac{5671}{35}\partial^5T+\hbar\,\mathrm{qt}(J^{(5)}),
\end{equation}
with generic-$c$ coefficients
$\beta_5^{(1)}=-c/56$, $\beta_5^{(2)}=-c(c+2)/280$.
\end{theorem}

\begin{proof}
The two-leg weight-$5$ space modulo total derivatives is empty;
the lowest-level nondegenerate weight-$5$ quasi-primary has
three $T$-legs. Solving $L_1e_5=0$ on the ansatz
${:}TT\partial T{:}+u{:}T\partial^3T{:}+v{:}(\partial T)(\partial^2T){:}
+w\partial^5T$:
\begin{itemize}
\item The ${:}T\Lambda_Z{:}$-coefficient closes iff $u=-3/10$
(the Zamolodchikov projection factor).
\item The ${:}(\partial T)(\partial^2T){:}$-sector receives
only the central contraction $T_{(3)}\partial^2T$, fixing
$v=-c/56$.
\item The $\partial^5T$-coefficient gathers two central tails,
closing on $w=-c(c+2)/280$.
\end{itemize}
At $c=-214$: $v=214/56=107/28$,
$w=-(-214)(-212)/280=-45368/280=-5671/35$.
Primary: Wang~$1998$~CMP~$195$ Prop.~$4.2$;
Bouwknegt--Schoutens~$1993$~Phys.\ Rep.~$223$.
\end{proof}

\begin{theorem}[Weight-$5$ pairing]
\label{thm:chi5-e5-vol3}
\ClaimStatusProvedHere
At $c=-214$,
\begin{equation}
\langle[\chi_5],[e_5]\rangle_{\Phi_5}
=\tfrac{5671}{35}\,\mathrm{Vol}(E)(2\pi\mathrm i)^5.
\end{equation}
\end{theorem}

\begin{proof}
The CoHA-Casimir at weight~$5$ equals
$\chi(\mathcal O_{\mathrm K3})\cdot|\beta_5^{(2)}|\cdot
\mathrm{Vol}(E)(2\pi\mathrm i)^5
=2\cdot 5671/70\cdot\mathrm{Vol}(E)(2\pi\mathrm i)^5
=5671/35\cdot\mathrm{Vol}(E)(2\pi\mathrm i)^5$.
Oberdieck's polar coefficient
$[p^{-1}q^{-1}y^5](-\Phi_{10}^{-1})$ returns the same rational
after Arnold normalisation. CY-$2$ formality
$\mathrm{Vol}(E)\to 0$ gives linear leading behaviour.
\end{proof}

\section{Weight-$6$ cocycle}
\label{sec:e6-calculus}

\begin{theorem}[$e_6$ chain-level form]
\label{thm:e6-vol3}
\ClaimStatusProvedHere
The weight-$6$ two-leg obstruction cocycle is
\begin{equation}
\label{eq:e6-vol3}
e_6(z)={:}T\partial^4T{:}-10\,{:}(\partial T)(\partial^3T){:}
     +10\,{:}(\partial^2T)^2{:}
     -\tfrac{7704}{7}\partial^6T+\hbar\,\mathrm{qt}(J^{(6)}),
\end{equation}
with generic-$c$ $\partial^6T$-coefficient $-c(c-2)/42$.
\end{theorem}

\begin{proof}
$L_1$ acts as
$L_1{:}T\partial^4T{:}=20{:}T\partial^3T{:}+\text{central}$,
$L_1{:}(\partial T)(\partial^3T){:}=2{:}T\partial^3T{:}+
12{:}(\partial T)(\partial^2T){:}+\text{central}$,
$L_1{:}(\partial^2T)^2{:}=12{:}(\partial T)(\partial^2T){:}+
\text{central}$, $L_1\partial^6T=42\partial^5T$.
Vanishing of ${:}T\partial^3T{:}$: $20+2\mu=0$, $\mu=-10$.
Vanishing of ${:}(\partial T)(\partial^2T){:}$:
$12\mu+12\nu=0$, $\nu=10$.
Vanishing of $\partial^5T$ with all central tails collected:
$\rho=-c(c-2)/42$.
At $c=-214$:
$\rho=-(-214)(-216)/42=-46224/42=-7704/7$.
Primary: Zamolodchikov~$1985$~TMF~$65$ \S$3$--$4$;
Nakayama~$2004$~IJMPA~$19$ App.~A.
\end{proof}

\section{Identification with $\mathcal W_\infty[c=-214]$ primaries}
\label{sec:winfinity-ident}

\begin{theorem}[$W_\infty$ identification]
\label{thm:ek-winfinity-vol3}
\ClaimStatusProvedHere
At $c=-214$, with $\Lambda_Z={:}TT{:}-\tfrac{3}{10}\partial^2T$ the
Zamolodchikov weight-$4$ primary and $W_k$ the
Pope--Romans--Shen~$1990$ integer-weight
$\mathcal W_\infty[c]$-primary,
\begin{enumerate}
\item $e_4=W_4-\tfrac{107}{11}\Lambda_Z$;
\item $e_5=W_5+\hbar\,\mathrm{qt}(J^{(5)})$;
\item $e_6=W_6-\tfrac{107}{11}\partial^2\Lambda_Z
         +\tfrac{23112}{(-1048)\cdot 42/5}{:}T\Lambda_Z{:}$,
with the denominator $c+22/5=-1048/5\ne 0$ regular at $c=-214$.
\end{enumerate}
\end{theorem}

\begin{proof}
Pope--Romans--Shen~$1990$~PLB~$236$ Eqs.~$(2.7$--$2.9)$ project
the quasi-primary onto the primary by subtracting
$\Lambda_Z$-descendants with scalar $-c/22$; at $c=-214$
this scalar is $214/22=107/11$. Three-leg weight-$5$
uniqueness (Wang~$1998$ Prop.~$4.2$) forces
$e_5=W_5$ identically.
\end{proof}

\begin{remark}[Two bases for $e_4$]
\label{rem:e4-two-bases}
Theorem~\ref{thm:e4-vol3} and item~(1) of
Theorem~\ref{thm:ek-winfinity-vol3} present the same class
$[e_4]\in\ChirHoch^3(\mathbf H_{\Delta_5})_4$ in two bases.
The Virasoro-composite basis
$({:}T\partial^2T{:},{:}(\partial T)^2{:},\partial^4T)$ carries
the quasi-primary coefficients $(1,-3/2,321/10)$, $L_1$-annihilated
by the computation of Theorem~\ref{thm:e4-vol3}. The primary basis
adds a fourth element $\Lambda_Z={:}TT{:}-(3/10)\partial^2T$, on
which the conversion reads
\[
W_4 \;=\; e_4 \;+\; (-c/22)\,\Lambda_Z
\qquad\text{(Pope--Romans--Shen $1990$).}
\]
At $c=-214$ the scalar is $214/22=107/11$, giving the advertised
$e_4=W_4-(107/11)\Lambda_Z$. Since $\Lambda_Z$ lies in the
$({:}TT{:},\partial^2T)$-sector, the $\partial^4T$-coefficient
$321/10$ is unshifted under the basis change; the $({:}T\partial^2T{:},{:}(\partial T)^2{:})$-coefficients $(1,-3/2)$
are likewise unshifted. The $W$-basis condenses the
$\Lambda_Z$-descendants into the single scalar $-c/22$;
the quasi-primary basis expands them explicitly. Both
expressions are the theorem.
\end{remark}

\section{Non-vanishing on K$3\times E$}
\label{sec:non-vanishing-vol3}

\begin{proposition}[Three-path non-vanishing of $[e_k]$ at $c=-214$]
\label{prop:ek-nonvanish}
\ClaimStatusProvedHere
For $k=4,5,6$, $[e_k]\ne 0$ in
$\ChirHoch^3(\mathbf H_{\Delta_5})_{k}$ at $c=-214$.
\end{proposition}

\begin{proof}
The pairings
$\langle[\chi_k],[e_k]\rangle_{\Phi_k}$ from
Theorems~\ref{thm:chi4-e4-vol3}, \ref{thm:chi5-e5-vol3}, and
(by the same argument at weight $6$)
$65193/10,5671/35,7704/7\cdot\mathrm{Vol}(E)(2\pi\mathrm i)^k$
are nonzero rationals times $\mathrm{Vol}(E)\ne 0$.
Non-vanishing of the pairing implies non-vanishing of the class.
\end{proof}

Primary:
Zamolodchikov~$1985$~TMF~$65$;
Bais--Bouwknegt--Schoutens--Surridge~$1988$~NPB~$304$;
Pope--Romans--Shen~$1990$~PLB~$236$;
Bouwknegt--Schoutens~$1993$~Phys.\ Rep.~$223$;
Wang~$1998$~CMP~$195$;
Nakayama~$2004$~IJMPA~$19$;
Schiffmann--Vasserot~$2013$~Publ.\ Math.\ IHES~$118$;
Oberdieck~$2018$~Geom.\ \& Topol.~$22$;
Kontsevich~$2003$~Lett.\ Math.\ Phys.~$66$.
